\section{Phương pháp Thực nghiệm}
\label{sec:experimental-methodology}

\subsection{Câu hỏi Nghiên cứu}
\label{sec:rqs}

\begin{itemize}
  \item \textbf{RQ1 (chất lượng).} Trên ISD-Agent-Bench, AlignmentGraph-ISD so với
  các ISD agent tiêu biểu về chất lượng artifact như thế nào, dưới cả một model lớn
  lẫn một small language model?
  \item \textbf{RQ2 (ablation).} Cơ chế nào của harness tạo ra chất lượng đó, và đóng
  góp của nó có thay đổi theo model scale không?
\end{itemize}

\subsection{Bộ dữ liệu}
\label{sec:dataset}

Chúng tôi dùng ISD-Agent-Bench~\cite{isdagentbench}, với các scenario trải một ma
trận Context~$\times$~ISD. Đánh giá chạy trên split \texttt{test\_90}---90 trong số
các test scenario, phân tầng theo 10 domain $\times$ 3 mức độ khó $\times$ 3 scenario
mỗi ô (seed 42)---cho phủ rộng, cân bằng với chi phí chấm hợp lý.

\subsection{Model và Thiết lập}
\label{sec:setup}

Để thăm dò tính độc lập với model, chúng tôi chạy mọi cấu hình ở hai điểm khác biệt
trên phổ model-scale: \textbf{DeepSeek-v4-flash} (một LLM lớn) và \textbf{Qwen3.5-9B}
(một SLM 9B), với cấu hình giống nhau trừ agent model. Agent model và judge model
được cấu hình riêng để tránh self-preference bias; các judge được giữ cố định qua mọi
lượt chạy và dùng cấu hình multi-judge mặc định của benchmark
(\texttt{openai/gpt-4o-mini} và \texttt{google/gemini-2.5-flash-lite}). Mỗi lượt chạy
ghi một manifest (model, judge, phiên bản prompt, commit, seed, cấu hình decoding)
để tái lập.

\subsection{Các Điều kiện}
\label{sec:conditions}

Với RQ1, AlignmentGraph-ISD được tích hợp làm agent thứ bảy dưới input/output schema
của benchmark và so với sáu agent: Baseline (single prompt), ReAct-ADDIE, ADDIE-Agent,
Dick-Carey-Agent, RPISD-Agent và EduPlanner.

Với RQ2, ablation loại bỏ đúng một cơ chế mỗi lần khỏi harness đầy đủ. Mỗi arm là một
workflow độc lập (không cờ cấu hình dùng chung), nên không cơ chế nào rò rỉ giữa các
arm:
\begin{itemize}
  \item \texttt{full} --- multi-agent decomposition + per-agent verify/repair;
  \item \texttt{wo-ma} --- bỏ decomposition: một agent đơn sinh toàn bộ chương trình
  và tự sửa toàn cục cho chính nó (giữ verify/repair);
  \item \texttt{wo-qc} --- bỏ verify/repair: cùng các specialist thứ tự cố định sinh,
  không có công cụ kiểm soát chất lượng.
\end{itemize}
So \texttt{full} với \texttt{wo-ma} tách riêng multi-agent decomposition; so với
\texttt{wo-qc} tách riêng per-agent verification/repair.

\subsection{Chỉ số}
\label{sec:metrics}

Các chỉ số chính là của benchmark: một ADDIE artifact rubric 100 điểm và một
trajectory score, gộp thành một \emph{Total} composite, mỗi cái là trung bình hai
judge. Chúng tôi bổ sung một \emph{residual alignment-defect audit} tất định:
verifier được chạy lại trên output cuối của mỗi arm và đếm các vi phạm mức nội dung
còn sót. Điều này đo trực tiếp điều kiện (i) của bài toán, nắm bắt một thuộc tính mà
rubric không tưởng thưởng tường minh.

\subsection{Phân tích Thống kê}
\label{sec:stats}

Các so sánh được ghép cặp theo scenario. Chúng tôi dùng kiểm định Wilcoxon
signed-rank (xấp xỉ chuẩn có hiệu chỉnh tie) trên composite score, với hiệu chỉnh
Holm qua các so sánh trong mỗi lượt chạy, và báo cáo hiệu trung bình theo cặp cùng
khoảng tin cậy bootstrap 95\% (10{,}000 resample). Với RQ2, đóng góp của một cơ chế là
$\texttt{full} - \texttt{arm}$, ghép cặp theo scenario.

\subsection{Tái lập}
\label{sec:reproducibility}

ISD-Agent-Bench và sáu baseline của nó là công khai. Mỗi lượt chạy được ghim bằng một
manifest (\cref{sec:setup}); ba ablation arm là các module workflow riêng chọn theo
tên; và các script phân tích (kiểm định ghép cặp và residual audit) được công bố cùng
bài báo.
